% Options for packages loaded elsewhere
\PassOptionsToPackage{unicode}{hyperref}
\PassOptionsToPackage{hyphens}{url}
%
\documentclass[
  10pt,
]{article}
\usepackage{amsmath,amssymb}
\usepackage{iftex}
\ifPDFTeX
  \usepackage[T1]{fontenc}
  \usepackage[utf8]{inputenc}
  \usepackage{textcomp} % provide euro and other symbols
\else % if luatex or xetex
  \usepackage{unicode-math} % this also loads fontspec
  \defaultfontfeatures{Scale=MatchLowercase}
  \defaultfontfeatures[\rmfamily]{Ligatures=TeX,Scale=1}
\fi
\ifPDFTeX\else
  % xetex/luatex font selection
  \setmainfont[]{DejaVu Serif}
  \setsansfont[]{DejaVu Sans}
  \setmonofont[]{DejaVu Sans Mono}
\fi
% Use upquote if available, for straight quotes in verbatim environments
\IfFileExists{upquote.sty}{\usepackage{upquote}}{}
\IfFileExists{microtype.sty}{% use microtype if available
  \usepackage[]{microtype}
  \UseMicrotypeSet[protrusion]{basicmath} % disable protrusion for tt fonts
}{}
\makeatletter
\@ifundefined{KOMAClassName}{% if non-KOMA class
  \IfFileExists{parskip.sty}{%
    \usepackage{parskip}
  }{% else
    \setlength{\parindent}{0pt}
    \setlength{\parskip}{6pt plus 2pt minus 1pt}}
}{% if KOMA class
  \KOMAoptions{parskip=half}}
\makeatother
\usepackage{xcolor}
\usepackage[margin=1in]{geometry}
\usepackage{color}
\usepackage{fancyvrb}
\newcommand{\VerbBar}{|}
\newcommand{\VERB}{\Verb[commandchars=\\\{\}]}
\DefineVerbatimEnvironment{Highlighting}{Verbatim}{commandchars=\\\{\}}
% Add ',fontsize=\small' for more characters per line
\newenvironment{Shaded}{}{}
\newcommand{\AlertTok}[1]{\textcolor[rgb]{1.00,0.00,0.00}{\textbf{#1}}}
\newcommand{\AnnotationTok}[1]{\textcolor[rgb]{0.38,0.63,0.69}{\textbf{\textit{#1}}}}
\newcommand{\AttributeTok}[1]{\textcolor[rgb]{0.49,0.56,0.16}{#1}}
\newcommand{\BaseNTok}[1]{\textcolor[rgb]{0.25,0.63,0.44}{#1}}
\newcommand{\BuiltInTok}[1]{\textcolor[rgb]{0.00,0.50,0.00}{#1}}
\newcommand{\CharTok}[1]{\textcolor[rgb]{0.25,0.44,0.63}{#1}}
\newcommand{\CommentTok}[1]{\textcolor[rgb]{0.38,0.63,0.69}{\textit{#1}}}
\newcommand{\CommentVarTok}[1]{\textcolor[rgb]{0.38,0.63,0.69}{\textbf{\textit{#1}}}}
\newcommand{\ConstantTok}[1]{\textcolor[rgb]{0.53,0.00,0.00}{#1}}
\newcommand{\ControlFlowTok}[1]{\textcolor[rgb]{0.00,0.44,0.13}{\textbf{#1}}}
\newcommand{\DataTypeTok}[1]{\textcolor[rgb]{0.56,0.13,0.00}{#1}}
\newcommand{\DecValTok}[1]{\textcolor[rgb]{0.25,0.63,0.44}{#1}}
\newcommand{\DocumentationTok}[1]{\textcolor[rgb]{0.73,0.13,0.13}{\textit{#1}}}
\newcommand{\ErrorTok}[1]{\textcolor[rgb]{1.00,0.00,0.00}{\textbf{#1}}}
\newcommand{\ExtensionTok}[1]{#1}
\newcommand{\FloatTok}[1]{\textcolor[rgb]{0.25,0.63,0.44}{#1}}
\newcommand{\FunctionTok}[1]{\textcolor[rgb]{0.02,0.16,0.49}{#1}}
\newcommand{\ImportTok}[1]{\textcolor[rgb]{0.00,0.50,0.00}{\textbf{#1}}}
\newcommand{\InformationTok}[1]{\textcolor[rgb]{0.38,0.63,0.69}{\textbf{\textit{#1}}}}
\newcommand{\KeywordTok}[1]{\textcolor[rgb]{0.00,0.44,0.13}{\textbf{#1}}}
\newcommand{\NormalTok}[1]{#1}
\newcommand{\OperatorTok}[1]{\textcolor[rgb]{0.40,0.40,0.40}{#1}}
\newcommand{\OtherTok}[1]{\textcolor[rgb]{0.00,0.44,0.13}{#1}}
\newcommand{\PreprocessorTok}[1]{\textcolor[rgb]{0.74,0.48,0.00}{#1}}
\newcommand{\RegionMarkerTok}[1]{#1}
\newcommand{\SpecialCharTok}[1]{\textcolor[rgb]{0.25,0.44,0.63}{#1}}
\newcommand{\SpecialStringTok}[1]{\textcolor[rgb]{0.73,0.40,0.53}{#1}}
\newcommand{\StringTok}[1]{\textcolor[rgb]{0.25,0.44,0.63}{#1}}
\newcommand{\VariableTok}[1]{\textcolor[rgb]{0.10,0.09,0.49}{#1}}
\newcommand{\VerbatimStringTok}[1]{\textcolor[rgb]{0.25,0.44,0.63}{#1}}
\newcommand{\WarningTok}[1]{\textcolor[rgb]{0.38,0.63,0.69}{\textbf{\textit{#1}}}}
\usepackage{longtable,booktabs,array}
\usepackage{calc} % for calculating minipage widths
% Correct order of tables after \paragraph or \subparagraph
\usepackage{etoolbox}
\makeatletter
\patchcmd\longtable{\par}{\if@noskipsec\mbox{}\fi\par}{}{}
\makeatother
% Allow footnotes in longtable head/foot
\IfFileExists{footnotehyper.sty}{\usepackage{footnotehyper}}{\usepackage{footnote}}
\makesavenoteenv{longtable}
\setlength{\emergencystretch}{3em} % prevent overfull lines
\providecommand{\tightlist}{%
  \setlength{\itemsep}{0pt}\setlength{\parskip}{0pt}}
\setcounter{secnumdepth}{-\maxdimen} % remove section numbering
\ifLuaTeX
  \usepackage{selnolig}  % disable illegal ligatures
\fi
\IfFileExists{bookmark.sty}{\usepackage{bookmark}}{\usepackage{hyperref}}
\IfFileExists{xurl.sty}{\usepackage{xurl}}{} % add URL line breaks if available
\urlstyle{same}
\hypersetup{
  hidelinks,
  pdfcreator={LaTeX via pandoc}}

\author{}
\date{}

\begin{document}

{
\setcounter{tocdepth}{3}
\tableofcontents
}
\hypertarget{dishtayantra-a-high-performance-multi-language-dataflow-framework-with-zero-copy-inter-process-communication}{%
\section{DishtaYantra: A High-Performance Multi-Language Dataflow
Framework with Zero-Copy Inter-Process
Communication}\label{dishtayantra-a-high-performance-multi-language-dataflow-framework-with-zero-copy-inter-process-communication}}

\textbf{Ashutosh Sinha}\\
Independent Researcher\\
ajsinha@gmail.com

\textbf{Repository:} https://github.com/ajsinha/dishtayantra

\begin{center}\rule{0.5\linewidth}{0.5pt}\end{center}

\hypertarget{abstract}{%
\subsection{Abstract}\label{abstract}}

Real-time data processing systems face fundamental challenges in
achieving low-latency computation while supporting heterogeneous
computational workloads across multiple programming languages. We
present \textbf{DishtaYantra}, a novel high-performance Directed Acyclic
Graph (DAG) compute framework that addresses these challenges through
three key innovations: (1) an LMDB-based zero-copy data exchange
mechanism achieving 100-1000× performance improvements over traditional
serialization for inter-language communication, (2) a unified
multi-language calculator framework supporting Python, Java (Py4J), C++
(pybind11), Rust (PyO3), and REST endpoints with sub-microsecond
invocation overhead, and (3) a multiprocessing worker pool architecture
with DAG affinity scheduling that bypasses Python's Global Interpreter
Lock (GIL) for true CPU parallelism. Our experimental evaluation on
consumer hardware (AMD Ryzen, 64GB RAM, Ubuntu 24.04) demonstrates that
DishtaYantra achieves median latencies of 5μs for large payload
transfers (compared to 500μs with JSON serialization), processes over
100,000 messages per second per worker, and maintains 90\% parallel
efficiency up to 12 concurrent workers. Around this compute core,
DishtaYantra provides a complete operational layer --- a pluggable
storage abstraction, database-backed authentication, configurable
high-availability with automatic failover, and market-aware scheduling
--- that allows the same dataflow definitions to run as a production
service. Two further advances extend this core: an \textbf{Arrow-native
columnar transport} that moves immutable record batches across DAG edges
\emph{by reference} --- preserving the unconditional change-only
recomputation gate while delivering a measured \textasciitilde2.3×
end-to-end speedup over the dictionary-envelope path with byte-identical
output --- and a \textbf{headless execution mode} with a
control-plane/worker orchestration layer that runs the same dataflow
definitions as bounded, idempotent batch jobs.

\textbf{Keywords:} Dataflow processing, DAG execution, Zero-copy
communication, Multi-language integration, LMDB, Real-time systems,
High-performance computing

\begin{center}\rule{0.5\linewidth}{0.5pt}\end{center}

\hypertarget{introduction}{%
\subsection{1. Introduction}\label{introduction}}

The proliferation of real-time data processing applications in financial
services, IoT analytics, and machine learning pipelines has created
unprecedented demand for high-performance dataflow frameworks {[}1,
2{]}. Modern enterprises require systems capable of processing millions
of events per second with sub-millisecond latencies while supporting
computational logic implemented in diverse programming languages
optimized for specific tasks {[}3{]}.

Existing dataflow frameworks face three fundamental limitations:

\begin{enumerate}
\def\labelenumi{\arabic{enumi}.}
\item
  \textbf{Serialization overhead} in inter-process and inter-language
  communication introduces latencies of 1-50ms for large payloads,
  making them unsuitable for microsecond-sensitive applications {[}4{]}.
\item
  \textbf{Language heterogeneity} forces developers to choose between
  optimal algorithms (e.g., numerical computing in C++/Rust) and
  framework integration (typically Python/Java), resulting in suboptimal
  performance or complex polyglot architectures {[}5{]}.
\item
  \textbf{GIL limitations} in Python-based frameworks prevent true CPU
  parallelism, constraining throughput regardless of available hardware
  resources {[}6{]}.
\end{enumerate}

We present \textbf{DishtaYantra}¹, a novel dataflow framework that
addresses these limitations through three key innovations:

\begin{enumerate}
\def\labelenumi{\arabic{enumi}.}
\item
  \textbf{Zero-Copy Data Exchange:} An LMDB-based transport mechanism
  that achieves 100-1000× performance improvements by eliminating
  serialization overhead through memory-mapped file I/O shared across
  process boundaries.
\item
  \textbf{Unified Multi-Language Support:} A calculator factory pattern
  supporting Python, Java (via Py4J), C++ (via pybind11), Rust (via
  PyO3), and external REST services with sub-100ns invocation overhead
  for native languages.
\item
  \textbf{GIL-Free Parallelism:} A multiprocessing worker pool with
  intelligent DAG affinity scheduling that enables true CPU parallelism
  by distributing workloads across independent Python interpreters.
\end{enumerate}

\begin{quote}
¹ Sanskrit for ``Computing Instrument'' -- reflecting the framework's
role as a precision tool for data processing.
\end{quote}

\hypertarget{contributions}{%
\subsubsection{1.1 Contributions}\label{contributions}}

This paper makes the following contributions:

\begin{itemize}
\tightlist
\item
  \textbf{DishtaYantra Framework:} A high-performance DAG compute
  framework supporting five programming languages with unified,
  format-agnostic configuration (§3)
\item
  \textbf{LMDB Zero-Copy Protocol:} A novel data exchange mechanism
  reducing latency by 100-1000× compared to JSON serialization (§6)
\item
  \textbf{Worker Pool Architecture:} DAG affinity scheduling with fault
  tolerance and automatic recovery (§7)
\item
  \textbf{Arrow-Native Columnar Transport:} A zero-copy, by-reference
  record-batch edge transport that preserves the change-only
  recomputation invariant while collapsing per-hop envelope overhead
  (§7B)
\item
  \textbf{Headless Execution \& Orchestration:} A webapp-free runner
  plus an idempotent, bounded control-plane/worker dispatch layer for
  batch and backfill workloads (§7B)
\item
  \textbf{Enterprise Operational Model:} A pluggable storage
  abstraction, database-backed authentication and authorization,
  configurable high-availability with automatic failover, and
  market-aware scheduling that together let the same dataflow run as a
  production service (§7A)
\item
  \textbf{Comprehensive Evaluation:} Performance characterization on
  consumer hardware (§9)
\end{itemize}

\begin{center}\rule{0.5\linewidth}{0.5pt}\end{center}

\hypertarget{background-and-motivation}{%
\subsection{2. Background and
Motivation}\label{background-and-motivation}}

\hypertarget{dataflow-processing-paradigm}{%
\subsubsection{2.1 Dataflow Processing
Paradigm}\label{dataflow-processing-paradigm}}

Dataflow processing systems model computation as directed acyclic graphs
(DAGs) where nodes represent computational operations and edges
represent data dependencies {[}7{]}. This paradigm offers several
advantages over imperative programming models:

\begin{itemize}
\tightlist
\item
  \textbf{Explicit Parallelism:} Independent nodes can execute
  concurrently without explicit synchronization
\item
  \textbf{Natural Pipeline Expression:} Data transformations are
  composed declaratively
\item
  \textbf{Clear Separation of Concerns:} Computation logic is decoupled
  from orchestration
\item
  \textbf{Fault Isolation:} Node failures can be isolated and recovered
  independently
\end{itemize}

\begin{verbatim}
                     Dataflow Processing Model
    
         [Source A] ──┐
                      ├──▶ [Transform] ──▶ [Aggregate] ──▶ [Sink]
         [Source B] ──┘         │
                                │
                                ▼
                           [Monitor]
\end{verbatim}

\hypertarget{the-serialization-bottleneck}{%
\subsubsection{2.2 The Serialization
Bottleneck}\label{the-serialization-bottleneck}}

Inter-process communication traditionally relies on serialization
formats such as JSON, Protocol Buffers, or Apache Avro {[}5{]}. While
these formats provide language interoperability, they introduce
significant overhead:

\begin{longtable}[]{@{}llll@{}}
\toprule\noalign{}
Operation & 1KB Payload & 100KB Payload & 1MB Payload \\
\midrule\noalign{}
\endhead
\bottomrule\noalign{}
\endlastfoot
JSON Serialize & 15 μs & 1.5 ms & 15 ms \\
JSON Deserialize & 20 μs & 2.0 ms & 20 ms \\
Network Transfer & 5 μs & 50 μs & 500 μs \\
\textbf{Total Overhead} & \textbf{40 μs} & \textbf{3.5 ms} & \textbf{35
ms} \\
\end{longtable}

For a 1MB payload, serialization and deserialization can consume 30-50ms
of wall-clock time, dominating the total latency budget in
microsecond-sensitive applications.

\hypertarget{python-gil-constraints}{%
\subsubsection{2.3 Python GIL
Constraints}\label{python-gil-constraints}}

Python's Global Interpreter Lock (GIL) is a mutex that protects access
to Python objects, preventing multiple native threads from executing
Python bytecode simultaneously {[}6{]}. This design choice simplifies
memory management but creates a fundamental throughput limitation:

\begin{verbatim}
                         Python GIL Limitation
    
    Thread 1: ████████░░░░░░░░████████░░░░░░░░████████░░░░
    Thread 2: ░░░░░░░░████████░░░░░░░░████████░░░░░░░░████
    Thread 3: ░░░░████░░░░████░░░░████░░░░████░░░░████░░░░
    
    ████ = Holding GIL (executing Python bytecode)
    ░░░░ = Blocked waiting for GIL
    
    Result: Only ONE thread executes Python at any given time
            regardless of available CPU cores
\end{verbatim}

DishtaYantra addresses this through multiprocessing, where each worker
runs in a separate Python interpreter with its own GIL.

\hypertarget{language-heterogeneity-challenge}{%
\subsubsection{2.4 Language Heterogeneity
Challenge}\label{language-heterogeneity-challenge}}

Modern data processing pipelines often require components written in
different languages:

\begin{itemize}
\tightlist
\item
  \textbf{Python:} Machine learning models (TensorFlow, PyTorch), data
  science (Pandas, NumPy)
\item
  \textbf{Java:} Enterprise systems, Kafka integration, existing
  business logic
\item
  \textbf{C++:} High-performance numerical computing, low-level
  optimizations, SIMD
\item
  \textbf{Rust:} Memory-safe systems programming, concurrent algorithms,
  zero-cost abstractions
\item
  \textbf{External Services:} REST APIs, microservices, legacy systems
\end{itemize}

Existing frameworks typically support only one or two languages, forcing
developers into suboptimal architectural choices.

\begin{center}\rule{0.5\linewidth}{0.5pt}\end{center}

\hypertarget{system-architecture}{%
\subsection{3. System Architecture}\label{system-architecture}}

DishtaYantra is architected as a layered system with clear separation of
concerns between presentation, orchestration, computation, and
infrastructure layers.

\hypertarget{high-level-architecture}{%
\subsubsection{3.1 High-Level
Architecture}\label{high-level-architecture}}

\begin{verbatim}
┌─────────────────────────────────────────────────────────────────────────────┐
│                           DishtaYantra v2.2                                │
├─────────────────────────────────────────────────────────────────────────────┤
│  PRESENTATION LAYER                                                          │
│  ┌──────────────┐  ┌──────────────┐  ┌──────────────┐  ┌─────────────────┐  │
│  │   Web UI     │  │  REST API    │  │   Admin      │  │   Prometheus    │  │
│  │  Dashboard   │  │  Endpoints   │  │  Dashboard   │  │   /metrics      │  │
│  └──────┬───────┘  └──────┬───────┘  └──────┬───────┘  └───────┬─────────┘  │
├─────────┴──────────────────┴──────────────────┴─────────────────┴───────────┤
│  APPLICATION LAYER                                                           │
│  ┌─────────────────────────────────────────────────────────────────────────┐ │
│  │  FastAPI Application │  Routes  │  Authentication  │  Session Mgmt      │ │
│  └─────────────────────────────────────────────────────────────────────────┘ │
├─────────────────────────────────────────────────────────────────────────────┤
│  ORCHESTRATION LAYER                                                         │
│  ┌───────────────┐  ┌───────────────┐  ┌───────────────┐  ┌──────────────┐  │
│  │  Worker Pool  │  │  DAG Affinity │  │    Health     │  │    Auto      │  │
│  │   Manager     │  │    Manager    │  │   Monitor     │  │   Recovery   │  │
│  └───────────────┘  └───────────────┘  └───────────────┘  └──────────────┘  │
├─────────────────────────────────────────────────────────────────────────────┤
│  COMPUTATION LAYER                                                           │
│  ┌───────────────┐  ┌───────────────┐  ┌───────────────┐  ┌──────────────┐  │
│  │ ComputeGraph  │  │    Node       │  │    Time       │  │   Subgraph   │  │
│  │   Engine      │  │   Executor    │  │   Windows     │  │   Support    │  │
│  └───────────────┘  └───────────────┘  └───────────────┘  └──────────────┘  │
├─────────────────────────────────────────────────────────────────────────────┤
│  INTEGRATION LAYER (Multi-Language Calculator Framework)                     │
│  ┌──────────┐ ┌──────────┐ ┌──────────┐ ┌──────────┐ ┌──────────┐          │
│  │  Python  │ │   Java   │ │   C++    │ │   Rust   │ │   REST   │          │
│  │ Built-in │ │  (Py4J)  │ │(pybind11)│ │  (PyO3)  │ │   API    │          │
│  └──────────┘ └──────────┘ └──────────┘ └──────────┘ └──────────┘          │
├─────────────────────────────────────────────────────────────────────────────┤
│  INFRASTRUCTURE LAYER                                                        │
│  ┌────────────────────────────────────────────────────────────────────────┐ │
│  │                    LMDB Zero-Copy Transport                             │ │
│  └────────────────────────────────────────────────────────────────────────┘ │
│  ┌───────┐ ┌───────┐ ┌───────┐ ┌───────┐ ┌───────┐ ┌───────┐ ┌──────────┐ │
│  │ Kafka │ │Rabbit │ │ Redis │ │Active │ │ TIBCO │ │IBM MQ │ │ InMemory │ │
│  │       │ │  MQ   │ │       │ │  MQ   │ │  EMS  │ │       │ │          │ │
│  └───────┘ └───────┘ └───────┘ └───────┘ └───────┘ └───────┘ └──────────┘ │
└─────────────────────────────────────────────────────────────────────────────┘
\end{verbatim}

\emph{Figure 1: DishtaYantra layered system architecture showing all six
architectural layers}

\hypertarget{core-abstractions}{%
\subsubsection{3.2 Core Abstractions}\label{core-abstractions}}

DishtaYantra is built around several fundamental abstractions that
provide a clean separation of concerns.

\hypertarget{computegraph}{%
\paragraph{3.2.1 ComputeGraph}\label{computegraph}}

The primary container representing a complete DAG:

\begin{verbatim}
┌─────────────────────────────────────────────────────────────────────┐
│                         ComputeGraph                                 │
├─────────────────────────────────────────────────────────────────────┤
│  Attributes:                                                         │
│  ┌─────────────────────────────────────────────────────────────────┐│
│  │ • name: str              - Unique identifier                    ││
│  │ • nodes: Dict[str, Node] - All nodes in the graph               ││
│  │ • edges: List[Edge]      - Data flow connections                ││
│  │ • status: GraphStatus    - STOPPED, RUNNING, SUSPENDED, ERROR   ││
│  │ • start_time: str        - Scheduled start (e.g., "0800")       ││
│  │ • duration: str          - Run duration (e.g., "12h")           ││
│  │ • auto_clone: bool       - Clone for next day automatically     ││
│  └─────────────────────────────────────────────────────────────────┘│
│                                                                      │
│  Methods:                                                            │
│  ┌─────────────────────────────────────────────────────────────────┐│
│  │ • start()     - Begin execution                                 ││
│  │ • stop()      - Graceful shutdown                               ││
│  │ • suspend()   - Pause without losing state                      ││
│  │ • resume()    - Continue from suspended state                   ││
│  │ • get_stats() - Execution statistics                            ││
│  └─────────────────────────────────────────────────────────────────┘│
└─────────────────────────────────────────────────────────────────────┘
\end{verbatim}

\hypertarget{node-types}{%
\paragraph{3.2.2 Node Types}\label{node-types}}

DishtaYantra provides six specialized node types:

\begin{verbatim}
┌─────────────────────────────────────────────────────────────────────┐
│                         Node Type Hierarchy                          │
├─────────────────────────────────────────────────────────────────────┤
│                                                                      │
│                            BaseNode                                  │
│                               │                                      │
│         ┌─────────┬──────────┼──────────┬──────────┬──────────┐    │
│         │         │          │          │          │          │     │
│         ▼         ▼          ▼          ▼          ▼          ▼     │
│  ┌───────────┐┌────────┐┌─────────┐┌─────────┐┌─────────┐┌───────┐ │
│  │Subscription││Publica-││Calcula- ││Transform││Metronome││Subgraph││
│  │   Node    ││  tion  ││  tion   ││  Node   ││  Node   ││ Node  │ │
│  │           ││  Node  ││  Node   ││         ││         ││       │ │
│  └───────────┘└────────┘└─────────┘└─────────┘└─────────┘└───────┘ │
│       │            │          │          │          │          │    │
│  ┌────┴────┐ ┌────┴────┐┌────┴────┐┌────┴────┐┌────┴────┐┌────┴──┐│
│  │Data     │ │Data     ││Calculator││Transform││Timer    ││Nested ││
│  │Ingestion│ │Output   ││Execution ││Logic    ││Events   ││DAG    ││
│  └─────────┘ └─────────┘└──────────┘└─────────┘└─────────┘└───────┘│
└─────────────────────────────────────────────────────────────────────┘
\end{verbatim}

\begin{longtable}[]{@{}
  >{\raggedright\arraybackslash}p{(\columnwidth - 4\tabcolsep) * \real{0.2821}}
  >{\raggedright\arraybackslash}p{(\columnwidth - 4\tabcolsep) * \real{0.2308}}
  >{\raggedright\arraybackslash}p{(\columnwidth - 4\tabcolsep) * \real{0.4872}}@{}}
\toprule\noalign{}
\begin{minipage}[b]{\linewidth}\raggedright
Node Type
\end{minipage} & \begin{minipage}[b]{\linewidth}\raggedright
Purpose
\end{minipage} & \begin{minipage}[b]{\linewidth}\raggedright
Key Configuration
\end{minipage} \\
\midrule\noalign{}
\endhead
\bottomrule\noalign{}
\endlastfoot
\textbf{SubscriptionNode} & Ingest data from external sources &
\texttt{subscribers:\ {[}\{uri:\ "kafka://..."\}{]}} \\
\textbf{PublicationNode} & Emit results to external sinks &
\texttt{publishers:\ {[}\{uri:\ "redis://..."\}{]}} \\
\textbf{CalculationNode} & Apply computational transformations &
\texttt{calculator:\ \{type:\ "rust",\ name:\ "..."\}} \\
\textbf{TransformNode} & Data structure transformations &
\texttt{transformer:\ \{field\_mappings:\ \{...\}\}} \\
\textbf{MetronomeNode} & Generate periodic trigger events &
\texttt{interval\_ms:\ 1000,\ payload:\ \{...\}} \\
\textbf{SubgraphNode} & Embed reusable sub-DAGs &
\texttt{subgraph\_path:\ "subgraphs/common.json"} \\
\end{longtable}

\hypertarget{edge-abstraction}{%
\paragraph{3.2.3 Edge Abstraction}\label{edge-abstraction}}

Edges define data flow with optional transformation:

\begin{verbatim}
┌─────────────────────────────────────────────────────────────────────┐
│                            Edge                                      │
├─────────────────────────────────────────────────────────────────────┤
│                                                                      │
│   Source Node                              Target Node               │
│  ┌──────────┐                             ┌──────────┐              │
│  │  Node A  │─────── Edge ───────────────▶│  Node B  │              │
│  └──────────┘                             └──────────┘              │
│                         │                                            │
│              ┌──────────┴──────────┐                                │
│              │  Edge Properties:    │                                │
│              │  • from: "node_a"    │                                │
│              │  • to: "node_b"      │                                │
│              │  • filter: {...}     │  (optional)                   │
│              │  • transform: {...}  │  (optional)                   │
│              │  • buffer_size: 1000 │  (flow control)               │
│              └─────────────────────┘                                │
└─────────────────────────────────────────────────────────────────────┘
\end{verbatim}

\hypertarget{calculator-interface}{%
\paragraph{3.2.4 Calculator Interface}\label{calculator-interface}}

All calculators implement a common interface:

\begin{Shaded}
\begin{Highlighting}[]
\KeywordTok{class}\NormalTok{ DataCalculator(ABC):}
    \CommentTok{"""Abstract base class for all calculators."""}
    
    \AttributeTok{@abstractmethod}
    \KeywordTok{def}\NormalTok{ calculate(}\VariableTok{self}\NormalTok{, data: }\BuiltInTok{dict}\NormalTok{) }\OperatorTok{{-}\textgreater{}} \BuiltInTok{dict}\NormalTok{:}
        \CommentTok{"""Transform input data and return result."""}
        \ControlFlowTok{pass}
    
    \AttributeTok{@abstractmethod}
    \KeywordTok{def}\NormalTok{ details(}\VariableTok{self}\NormalTok{) }\OperatorTok{{-}\textgreater{}} \BuiltInTok{dict}\NormalTok{:}
        \CommentTok{"""Return calculator metadata and runtime statistics."""}
        \ControlFlowTok{pass}
\end{Highlighting}
\end{Shaded}

\hypertarget{pubsub-framework}{%
\subsubsection{3.3 Pub/Sub Framework}\label{pubsub-framework}}

DishtaYantra provides a unified publish/subscribe abstraction supporting
multiple message brokers:

\begin{verbatim}
┌─────────────────────────────────────────────────────────────────────┐
│                      Pub/Sub Framework Architecture                  │
├─────────────────────────────────────────────────────────────────────┤
│                                                                      │
│  URI Format: protocol://host:port/topic?params                       │
│                                                                      │
│  ┌─────────────────────────────────────────────────────────────────┐│
│  │                    DataPublisher / DataSubscriber               ││
│  │                         (Abstract Interface)                    ││
│  └───────────────────────────────┬─────────────────────────────────┘│
│                                  │                                   │
│    ┌────────┬────────┬──────────┼──────────┬────────┬────────┐     │
│    ▼        ▼        ▼          ▼          ▼        ▼        ▼      │
│ ┌──────┐┌──────┐┌──────┐┌──────────┐┌──────┐┌──────┐┌──────────┐   │
│ │Kafka ││Rabbit││Redis ││ ActiveMQ ││TIBCO ││IBM MQ││ InMemory │   │
│ └──────┘└──────┘└──────┘└──────────┘└──────┘└──────┘└──────────┘   │
│                                                                      │
│  ┌─────────────────────────────────────────────────────────────────┐│
│  │                    DataPublisher / DataSubscriber               ││
│  │                         (Abstract Interface)                    ││
│  └───────────────────────────────┬─────────────────────────────────┘│
│                                  │                                   │
│  Brokers:  kafka  rabbitmq  activemq  tibcoems  websphere  redis     │
│  Cloud:    AWS (sqs, kinesis, sns)   Azure (servicebus, eventhubs)   │
│  Stores:   s3  azureblob  gcs   Other: file, sql, aerospike, grpc    │
│  Local:    inmemory/mem, lmdb (zero-copy)   Compose: fanin, custom   │
│  Wrapping: any of the above with "resilient": true                   │
└─────────────────────────────────────────────────────────────────────┘
\end{verbatim}

\textbf{Representative backends} (latencies are indicative, consumer
hardware):

\begin{longtable}[]{@{}
  >{\raggedright\arraybackslash}p{(\columnwidth - 6\tabcolsep) * \real{0.2632}}
  >{\raggedright\arraybackslash}p{(\columnwidth - 6\tabcolsep) * \real{0.2368}}
  >{\raggedright\arraybackslash}p{(\columnwidth - 6\tabcolsep) * \real{0.2632}}
  >{\raggedright\arraybackslash}p{(\columnwidth - 6\tabcolsep) * \real{0.2368}}@{}}
\toprule\noalign{}
\begin{minipage}[b]{\linewidth}\raggedright
Protocol
\end{minipage} & \begin{minipage}[b]{\linewidth}\raggedright
Backend
\end{minipage} & \begin{minipage}[b]{\linewidth}\raggedright
Use Case
\end{minipage} & \begin{minipage}[b]{\linewidth}\raggedright
Latency
\end{minipage} \\
\midrule\noalign{}
\endhead
\bottomrule\noalign{}
\endlastfoot
\texttt{kafka://} & Apache Kafka & High-throughput streaming &
\textasciitilde5ms \\
\texttt{rabbitmq://} & RabbitMQ & Enterprise messaging &
\textasciitilde1ms \\
\texttt{activemq://} & ActiveMQ & JMS compatibility &
\textasciitilde2ms \\
\texttt{tibcoems://} & TIBCO EMS & Enterprise integration &
\textasciitilde1ms \\
\texttt{websphere://} & IBM MQ & Mainframe connectivity &
\textasciitilde5ms \\
\texttt{redis://}, \texttt{redischannel://} & Redis & Low-latency
caching/pub-sub & \textasciitilde100μs \\
\texttt{sqs://}, \texttt{kinesis://}, \texttt{sns://} & AWS managed
messaging & Cloud-native streaming/queues & network-bound \\
\texttt{servicebus://}, \texttt{eventhubs://} & Azure managed messaging
& Cloud-native messaging & network-bound \\
\texttt{s3://}, \texttt{azureblob://}, \texttt{gcs://} & Cloud object
stores & Durable polling pub/sub & network-bound \\
\texttt{sql://}, \texttt{aerospike://}, \texttt{grpc://} & Database / KV
/ RPC & Integration sources/sinks & varies \\
\texttt{file://} & Filesystem & Durable local I/O & disk-bound \\
\texttt{lmdb://} & LMDB & Cross-process IPC (zero-copy) &
\textasciitilde5μs \\
\texttt{inmemory://}, \texttt{mem://} & In-Memory & Testing/Development
& \textasciitilde1μs \\
\texttt{fanin://} & Composite & Merge several subscribers & n/a \\
\end{longtable}

Any broker/cloud backend can be wrapped with a \textbf{resilient}
variant (\texttt{"resilient":\ true}) that adds automatic reconnection,
message buffering during outages, and subscription/state restoration on
reconnect.

\hypertarget{backpressure-management}{%
\subsubsection{3.4 Backpressure
Management}\label{backpressure-management}}

DishtaYantra implements comprehensive backpressure to prevent memory
exhaustion:

\begin{verbatim}
┌─────────────────────────────────────────────────────────────────────┐
│                    Backpressure Management                           │
├─────────────────────────────────────────────────────────────────────┤
│                                                                      │
│  Producer                 Queue                   Consumer           │
│  ┌──────┐            ┌──────────────┐            ┌──────┐           │
│  │      │ ──────────▶│ ████████░░░░ │──────────▶ │      │           │
│  │ Fast │            │ (80% full)   │            │ Slow │           │
│  │      │ ◀── SLOW ──│              │            │      │           │
│  └──────┘    DOWN    └──────────────┘            └──────┘           │
│                             │                                        │
│                    Backpressure Signal                               │
│                                                                      │
│  Strategies:                                                         │
│  ┌─────────────────────────────────────────────────────────────────┐│
│  │ • BLOCK:    Upstream blocks when queue full (default)           ││
│  │ • DROP:     Excess messages discarded with logging              ││
│  │ • SAMPLE:   Statistical sampling under high load                ││
│  │ • THROTTLE: Rate limiting at source nodes                       ││
│  └─────────────────────────────────────────────────────────────────┘│
└─────────────────────────────────────────────────────────────────────┘
\end{verbatim}

\hypertarget{internal-cache-for-crash-recovery}{%
\subsubsection{3.5 Internal Cache for Crash
Recovery}\label{internal-cache-for-crash-recovery}}

DishtaYantra maintains an internal cache enabling recovery without data
loss:

\begin{verbatim}
┌─────────────────────────────────────────────────────────────────────┐
│                    Crash Recovery Architecture                       │
├─────────────────────────────────────────────────────────────────────┤
│                                                                      │
│  Normal Operation:                                                   │
│  ┌────────────┐      ┌────────────┐      ┌────────────┐            │
│  │  Process   │─────▶│   Cache    │─────▶│  Persist   │            │
│  │   Data     │      │  In-Flight │      │ Checkpoint │            │
│  └────────────┘      └────────────┘      └────────────┘            │
│                                                                      │
│  Cache Contents:                                                     │
│  ┌─────────────────────────────────────────────────────────────────┐│
│  │ • In-flight messages: Data currently being processed            ││
│  │ • Node state: Calculator state for stateful computations        ││
│  │ • Processing offsets: Position markers for replay               ││
│  │ • Checkpoint timestamp: Last successful commit                  ││
│  └─────────────────────────────────────────────────────────────────┘│
│                                                                      │
│  Recovery Sequence:                                                  │
│  1. Load last checkpoint from cache                                  │
│  2. Identify in-flight messages                                      │
│  3. Replay from recovery point                                       │
│  4. Resume normal processing (at-least-once semantics)               │
│                                                                      │
│  Backend Options: LMDB (default), Redis, In-Memory                   │
└─────────────────────────────────────────────────────────────────────┘
\end{verbatim}

\begin{center}\rule{0.5\linewidth}{0.5pt}\end{center}

\hypertarget{dag-definition-and-execution}{%
\subsection{4. DAG Definition and
Execution}\label{dag-definition-and-execution}}

\hypertarget{json-configuration-format}{%
\subsubsection{4.1 JSON Configuration
Format}\label{json-configuration-format}}

DishtaYantra DAGs are defined declaratively using JSON:

\begin{Shaded}
\begin{Highlighting}[]
\FunctionTok{\{}
    \DataTypeTok{"name"}\FunctionTok{:} \StringTok{"trade\_pricing\_pipeline"}\FunctionTok{,}
    \DataTypeTok{"description"}\FunctionTok{:} \StringTok{"Real{-}time trade pricing with risk calculation"}\FunctionTok{,}
    \DataTypeTok{"start\_time"}\FunctionTok{:} \StringTok{"0800"}\FunctionTok{,}
    \DataTypeTok{"duration"}\FunctionTok{:} \StringTok{"12h"}\FunctionTok{,}
    \DataTypeTok{"nodes"}\FunctionTok{:} \OtherTok{[}
        \FunctionTok{\{}
            \DataTypeTok{"name"}\FunctionTok{:} \StringTok{"trade\_subscriber"}\FunctionTok{,}
            \DataTypeTok{"type"}\FunctionTok{:} \StringTok{"SubscriptionNode"}\FunctionTok{,}
            \DataTypeTok{"subscribers"}\FunctionTok{:} \OtherTok{[}\FunctionTok{\{}\DataTypeTok{"uri"}\FunctionTok{:} \StringTok{"kafka://broker:9092/raw\_trades"}\FunctionTok{\}}\OtherTok{]}
        \FunctionTok{\}}\OtherTok{,}
        \FunctionTok{\{}
            \DataTypeTok{"name"}\FunctionTok{:} \StringTok{"price\_calculator"}\FunctionTok{,}
            \DataTypeTok{"type"}\FunctionTok{:} \StringTok{"CalculationNode"}\FunctionTok{,}
            \DataTypeTok{"calculator"}\FunctionTok{:} \FunctionTok{\{}\DataTypeTok{"type"}\FunctionTok{:} \StringTok{"rust"}\FunctionTok{,} \DataTypeTok{"name"}\FunctionTok{:} \StringTok{"RustTradePricingCalculator"}\FunctionTok{\}}
        \FunctionTok{\}}\OtherTok{,}
        \FunctionTok{\{}
            \DataTypeTok{"name"}\FunctionTok{:} \StringTok{"risk\_calculator"}\FunctionTok{,}
            \DataTypeTok{"type"}\FunctionTok{:} \StringTok{"CalculationNode"}\FunctionTok{,}
            \DataTypeTok{"calculator"}\FunctionTok{:} \FunctionTok{\{}\DataTypeTok{"type"}\FunctionTok{:} \StringTok{"cpp"}\FunctionTok{,} \DataTypeTok{"name"}\FunctionTok{:} \StringTok{"CppRiskMetricsCalculator"}\FunctionTok{\}}
        \FunctionTok{\}}\OtherTok{,}
        \FunctionTok{\{}
            \DataTypeTok{"name"}\FunctionTok{:} \StringTok{"result\_publisher"}\FunctionTok{,}
            \DataTypeTok{"type"}\FunctionTok{:} \StringTok{"PublicationNode"}\FunctionTok{,}
            \DataTypeTok{"publishers"}\FunctionTok{:} \OtherTok{[}
                \FunctionTok{\{}\DataTypeTok{"uri"}\FunctionTok{:} \StringTok{"kafka://broker:9092/priced\_trades"}\FunctionTok{\}}\OtherTok{,}
                \FunctionTok{\{}\DataTypeTok{"uri"}\FunctionTok{:} \StringTok{"redis://cache:6379/latest\_prices"}\FunctionTok{\}}
            \OtherTok{]}
        \FunctionTok{\}}
    \OtherTok{]}\FunctionTok{,}
    \DataTypeTok{"edges"}\FunctionTok{:} \OtherTok{[}
        \FunctionTok{\{}\DataTypeTok{"from"}\FunctionTok{:} \StringTok{"trade\_subscriber"}\FunctionTok{,} \DataTypeTok{"to"}\FunctionTok{:} \StringTok{"price\_calculator"}\FunctionTok{\}}\OtherTok{,}
        \FunctionTok{\{}\DataTypeTok{"from"}\FunctionTok{:} \StringTok{"price\_calculator"}\FunctionTok{,} \DataTypeTok{"to"}\FunctionTok{:} \StringTok{"risk\_calculator"}\FunctionTok{\}}\OtherTok{,}
        \FunctionTok{\{}\DataTypeTok{"from"}\FunctionTok{:} \StringTok{"risk\_calculator"}\FunctionTok{,} \DataTypeTok{"to"}\FunctionTok{:} \StringTok{"result\_publisher"}\FunctionTok{\}}
    \OtherTok{]}
\FunctionTok{\}}
\end{Highlighting}
\end{Shaded}

\hypertarget{visual-dag-representation}{%
\subsubsection{4.2 Visual DAG
Representation}\label{visual-dag-representation}}

\begin{verbatim}
┌─────────────────────────────────────────────────────────────────┐
│                    trade_pricing_pipeline DAG                    │
├─────────────────────────────────────────────────────────────────┤
│                                                                  │
│     ┌───────────────┐                                           │
│     │  Kafka Topic  │  (External Source)                        │
│     │  raw_trades   │                                           │
│     └───────┬───────┘                                           │
│             │                                                    │
│             ▼                                                    │
│     ┌───────────────────┐                                       │
│     │ trade_subscriber  │  SubscriptionNode                     │
│     └─────────┬─────────┘                                       │
│               │ trade data (JSON)                                │
│               ▼                                                  │
│     ┌───────────────────┐                                       │
│     │ price_calculator  │  CalculationNode (Rust)               │
│     └─────────┬─────────┘                                       │
│               │ priced trade                                     │
│               ▼                                                  │
│     ┌───────────────────┐                                       │
│     │ risk_calculator   │  CalculationNode (C++)                │
│     └─────────┬─────────┘                                       │
│               │ trade + risk metrics                             │
│               ▼                                                  │
│     ┌───────────────────┐                                       │
│     │ result_publisher  │  PublicationNode                      │
│     └─────────┬─────────┘                                       │
│               │                                                  │
│      ┌────────┴────────┐                                        │
│      ▼                 ▼                                         │
│ ┌──────────┐    ┌───────────┐                                   │
│ │  Kafka   │    │   Redis   │  (External Sinks)                 │
│ │ priced_  │    │  latest_  │                                   │
│ │ trades   │    │  prices   │                                   │
│ └──────────┘    └───────────┘                                   │
└─────────────────────────────────────────────────────────────────┘
\end{verbatim}

\emph{Figure 2: Trade pricing pipeline DAG showing data flow from Kafka
source through Rust pricing and C++ risk calculators to multiple output
destinations}

\hypertarget{dag-execution-algorithm}{%
\subsubsection{4.3 DAG Execution
Algorithm}\label{dag-execution-algorithm}}

The execution engine uses Kahn's algorithm for topological sorting:

\begin{verbatim}
Algorithm 1: DAG Execution Loop
─────────────────────────────────────────────────────────────────────
Input:  ComputeGraph G with nodes N and edges E
Output: Continuous processing of data through the graph

 1:  execution_order ← TopologicalSort(G)
 2:  for each node n in N do
 3:      n.dirty ← false
 4:      n.input_queue ← Queue(max_size=BUFFER_SIZE)
 5:  end for
 6:  
 7:  while G.status = RUNNING do
 8:      for each node n in execution_order do
 9:          if n.dirty OR n.has_new_input() then
10:              input_data ← n.collect_inputs()
11:              output_data ← n.process(input_data)
12:              n.output ← output_data
13:              n.dirty ← false
14:              for each downstream d in G.successors(n) do
15:                  d.dirty ← true
16:                  d.input_queue.push(output_data)
17:              end for
18:          end if
19:      end for
20:      sleep(G.poll_interval)
21:  end while
─────────────────────────────────────────────────────────────────────
\end{verbatim}

\begin{center}\rule{0.5\linewidth}{0.5pt}\end{center}

\hypertarget{multi-language-calculator-framework}{%
\subsection{5. Multi-Language Calculator
Framework}\label{multi-language-calculator-framework}}

\hypertarget{calculator-factory-architecture}{%
\subsubsection{5.1 Calculator Factory
Architecture}\label{calculator-factory-architecture}}

\begin{verbatim}
┌─────────────────────────────────────────────────────────────────────┐
│                       CalculatorFactory                              │
├─────────────────────────────────────────────────────────────────────┤
│                                                                      │
│  create(name, config) → DataCalculator                               │
│                                                                      │
│              config.calculator.type                                  │
│                       │                                              │
│    ┌──────┬──────┬───┴───┬──────┬──────┐                            │
│    ▼      ▼      ▼       ▼      ▼      ▼                             │
│  "java" "cpp" "rust"  "rest" "python" (default)                      │
│    │      │      │       │      │      │                             │
│    ▼      ▼      ▼       ▼      ▼      ▼                             │
│  Java   Cpp   Rust    Rest  Built-in Built-in                        │
│  Calc   Calc  Calc    Calc  Calculator                               │
│ (Py4J) (pybind)(PyO3) (HTTP)                                         │
└─────────────────────────────────────────────────────────────────────┘
\end{verbatim}

\hypertarget{java-integration-py4j}{%
\subsubsection{5.2 Java Integration
(Py4J)}\label{java-integration-py4j}}

\begin{verbatim}
┌─────────────────────────────────────────────────────────────────────┐
│                     Java Integration via Py4J                        │
├─────────────────────────────────────────────────────────────────────┤
│                                                                      │
│  Python Process                    │    JVM Process                  │
│  ┌─────────────────────────┐       │    ┌─────────────────────────┐ │
│  │     JavaCalculator      │       │    │  DishtaYantraGateway    │ │
│  │  ┌───────────────────┐  │ Py4J  │    │  ┌───────────────────┐  │ │
│  │  │  gateway_pool     │◀─┼───────┼───▶│  │ Calculator        │  │ │
│  │  │  (connection pool)│  │ TCP   │    │  │ Registry          │  │ │
│  │  └───────────────────┘  │ 25333 │    │  └───────────────────┘  │ │
│  └─────────────────────────┘       │    └─────────────────────────┘ │
│                                    │                                 │
│  Performance: ~1ms per call        │    Heap: 2GB (configurable)    │
│  Connection pool size: 4           │    GC: G1GC recommended        │
└─────────────────────────────────────────────────────────────────────┘
\end{verbatim}

\hypertarget{c-integration-pybind11}{%
\subsubsection{5.3 C++ Integration
(pybind11)}\label{c-integration-pybind11}}

\begin{verbatim}
┌─────────────────────────────────────────────────────────────────────┐
│                   C++ Integration via pybind11                       │
├─────────────────────────────────────────────────────────────────────┤
│                                                                      │
│  Python Process                                                      │
│  ┌─────────────────────────────────────────────────────────────────┐│
│  │  import dishtayantra_cpp                                        ││
│  │                                                                  ││
│  │  ┌────────────────────┐      ┌────────────────────────────────┐││
│  │  │    CppCalculator   │      │    dishtayantra_cpp.so         │││
│  │  │  • calculate()     │─────▶│    • MathCalculator            │││
│  │  │  • details()       │      │    • StatisticsCalculator      │││
│  │  └────────────────────┘      │    • TradePricingCalculator    │││
│  │                              │    • RiskMetricsCalculator     │││
│  │  Performance:                └────────────────────────────────┘││
│  │  • ~100ns call overhead                                         ││
│  │  • Direct memory access                                         ││
│  └─────────────────────────────────────────────────────────────────┘│
└─────────────────────────────────────────────────────────────────────┘
\end{verbatim}

\hypertarget{rust-integration-pyo3}{%
\subsubsection{5.4 Rust Integration
(PyO3)}\label{rust-integration-pyo3}}

\begin{verbatim}
┌─────────────────────────────────────────────────────────────────────┐
│                    Rust Integration via PyO3                         │
├─────────────────────────────────────────────────────────────────────┤
│                                                                      │
│  Python Process                                                      │
│  ┌─────────────────────────────────────────────────────────────────┐│
│  │  import dishtayantra_rust                                       ││
│  │                                                                  ││
│  │  ┌────────────────────┐      ┌────────────────────────────────┐││
│  │  │   RustCalculator   │      │  dishtayantra_rust.so          │││
│  │  │  • calculate()     │─────▶│    • MathCalculator            │││
│  │  │  • details()       │      │    • StatisticsCalculator      │││
│  │  └────────────────────┘      │    • TradePricingCalculator    │││
│  │                              └────────────────────────────────┘││
│  │  Benefits: Memory safety, Thread safety, Zero-cost abstractions ││
│  └─────────────────────────────────────────────────────────────────┘│
└─────────────────────────────────────────────────────────────────────┘
\end{verbatim}

\hypertarget{rest-integration}{%
\subsubsection{5.5 REST Integration}\label{rest-integration}}

\begin{verbatim}
┌─────────────────────────────────────────────────────────────────────┐
│                     REST Integration via HTTP                        │
├─────────────────────────────────────────────────────────────────────┤
│                                                                      │
│  DishtaYantra                         External Service               │
│  ┌─────────────────────────┐          ┌─────────────────────────┐   │
│  │     RestCalculator      │   HTTP   │    REST Endpoint        │   │
│  │  endpoint: /api/calc    │   POST   │    POST /api/calc       │   │
│  │  timeout: 30s           │ ────────▶│    - Process request    │   │
│  │  auth_type: bearer      │ ◀────────│    - Return result      │   │
│  └─────────────────────────┘   JSON   └─────────────────────────┘   │
│                                                                      │
│  Authentication: api_key, basic, bearer, custom                      │
└─────────────────────────────────────────────────────────────────────┘
\end{verbatim}

\begin{center}\rule{0.5\linewidth}{0.5pt}\end{center}

\hypertarget{lmdb-zero-copy-data-exchange}{%
\subsection{6. LMDB Zero-Copy Data
Exchange}\label{lmdb-zero-copy-data-exchange}}

\hypertarget{design-principles}{%
\subsubsection{6.1 Design Principles}\label{design-principles}}

Traditional inter-process communication requires serialization at the
sender and deserialization at the receiver. DishtaYantra's LMDB-based
approach eliminates this overhead by leveraging memory-mapped files.

\textbf{Key insight:} LMDB {[}8{]} provides a lightning-fast,
memory-mapped key-value store with ACID guarantees. By writing data once
to LMDB and passing only the key reference to native calculators, we
achieve zero-copy reads across process and language boundaries.

\hypertarget{architecture}{%
\subsubsection{6.2 Architecture}\label{architecture}}

\begin{verbatim}
┌─────────────────────────────────────────────────────────────────────┐
│                  LMDB Zero-Copy Data Exchange                        │
├─────────────────────────────────────────────────────────────────────┤
│                                                                      │
│  Step 1: Python writes data to LMDB                                  │
│  ┌─────────────────────────────────────────────────────────────────┐│
│  │  Python DAG Engine                                               ││
│  │  ┌────────────┐         ┌────────────────────────────┐          ││
│  │  │ Input Data │────────▶│     LMDBTransport          │          ││
│  │  │ (dict/list)│  write  │     serialize + store      │          ││
│  │  └────────────┘         └─────────────┬──────────────┘          ││
│  └───────────────────────────────────────┼──────────────────────────┘│
│                                          ▼                           │
│  ┌─────────────────────────────────────────────────────────────────┐│
│  │                    LMDB File (Memory-Mapped)                     ││
│  │  Key: "calc_12345_input"  │  Value: <binary data>                ││
│  │              │                                                   ││
│  │              │  Memory Map (mmap) - Zero-copy access             ││
│  │              ▼                                                   ││
│  └─────────────────────────────────────────────────────────────────┘│
│                                                                      │
│  Step 2: Native calculator reads directly from mapped memory         │
│  ┌─────────────────────────────────────────────────────────────────┐│
│  │  ┌────────────┐    ┌────────────┐    ┌────────────────────┐    ││
│  │  │    Java    │    │    C++     │    │       Rust         │    ││
│  │  │  (lmdbjava)│    │  (liblmdb) │    │    (lmdb-rs)       │    ││
│  │  │  Direct    │    │  Direct    │    │  Direct memory     │    ││
│  │  │  memory    │    │  memory    │    │  access via mmap   │    ││
│  │  └────────────┘    └────────────┘    └────────────────────┘    ││
│  └─────────────────────────────────────────────────────────────────┘│
│                                                                      │
│  Step 3: Output written back to LMDB, Python reads result            │
└─────────────────────────────────────────────────────────────────────┘
\end{verbatim}

\emph{Figure 3: LMDB zero-copy data exchange architecture}

\hypertarget{exchange-modes}{%
\subsubsection{6.3 Exchange Modes}\label{exchange-modes}}

\begin{longtable}[]{@{}llll@{}}
\toprule\noalign{}
Mode & Input & Output & Use Case \\
\midrule\noalign{}
\endhead
\bottomrule\noalign{}
\endlastfoot
\texttt{input} & LMDB & Direct & Large inputs, small outputs \\
\texttt{output} & Direct & LMDB & Small inputs, large outputs \\
\texttt{both} & LMDB & LMDB & Large data in both directions \\
\texttt{reference} & Key only & Key only & Calculator manages all I/O \\
\end{longtable}

\hypertarget{performance-results}{%
\subsubsection{6.4 Performance Results}\label{performance-results}}

\begin{longtable}[]{@{}llll@{}}
\toprule\noalign{}
Payload Size & JSON Serialization & LMDB Zero-Copy & Speedup \\
\midrule\noalign{}
\endhead
\bottomrule\noalign{}
\endlastfoot
1 KB & 50 μs & \textbf{5 μs} & 10× \\
10 KB & 500 μs & \textbf{10 μs} & 50× \\
100 KB & 5 ms & \textbf{50 μs} & 100× \\
1 MB & 50 ms & \textbf{200 μs} & 250× \\
10 MB & 500 ms & \textbf{2 ms} & 250× \\
\end{longtable}

\hypertarget{configuration}{%
\subsubsection{6.5 Configuration}\label{configuration}}

\begin{Shaded}
\begin{Highlighting}[]
\NormalTok{\# application.properties}
\NormalTok{lmdb.db.path=/tmp/dishtayantra\_lmdb}
\NormalTok{lmdb.map.size=1073741824          \# 1GB}
\NormalTok{lmdb.max.dbs=100}
\NormalTok{lmdb.ttl.seconds=300}
\NormalTok{lmdb.cleanup.interval=60}
\end{Highlighting}
\end{Shaded}

\begin{center}\rule{0.5\linewidth}{0.5pt}\end{center}

\hypertarget{worker-pool-architecture}{%
\subsection{7. Worker Pool
Architecture}\label{worker-pool-architecture}}

\hypertarget{overview}{%
\subsubsection{7.1 Overview}\label{overview}}

DishtaYantra's worker pool enables true CPU parallelism by running DAGs
in separate Python interpreter processes:

\begin{verbatim}
┌─────────────────────────────────────────────────────────────────────┐
│                    Worker Pool Architecture                          │
├─────────────────────────────────────────────────────────────────────┤
│                                                                      │
│  Main Process                                                        │
│  ┌─────────────────────────────────────────────────────────────────┐│
│  │  WorkerPoolManager                                               ││
│  │  ┌─────────────┐  ┌─────────────┐  ┌─────────────────────────┐ ││
│  │  │  Affinity   │  │   Health    │  │     Auto-Recovery       │ ││
│  │  │  Manager    │  │   Monitor   │  │     (restart failed)    │ ││
│  │  └─────────────┘  └─────────────┘  └─────────────────────────┘ ││
│  └──────────────────────────┬──────────────────────────────────────┘│
│                             │                                        │
│        ┌────────────────────┼────────────────────┐                  │
│        ▼                    ▼                    ▼                   │
│  ┌──────────────┐    ┌──────────────┐    ┌──────────────┐          │
│  │   Worker 0   │    │   Worker 1   │    │   Worker N   │          │
│  │  (Process)   │    │  (Process)   │    │  (Process)   │          │
│  │  Own GIL     │    │  Own GIL     │    │  Own GIL     │          │
│  │  Own Memory  │    │  Own Memory  │    │  Own Memory  │          │
│  │  ┌────────┐  │    │  ┌────────┐  │    │  ┌────────┐  │          │
│  │  │ DAG A  │  │    │  │ DAG C  │  │    │  │ DAG E  │  │          │
│  │  │ DAG B  │  │    │  │ DAG D  │  │    │  │ DAG F  │  │          │
│  │  └────────┘  │    │  └────────┘  │    │  └────────┘  │          │
│  └──────────────┘    └──────────────┘    └──────────────┘          │
│         │                   │                   │                   │
│         └───────────────────┼───────────────────┘                   │
│                    ┌────────┴────────┐                              │
│                    │  LMDB Pub/Sub   │                              │
│                    │ Cross-Worker IPC│                              │
│                    └─────────────────┘                              │
└─────────────────────────────────────────────────────────────────────┘
\end{verbatim}

\emph{Figure 4: Worker pool architecture with independent GILs}

\hypertarget{dag-affinity-scheduling}{%
\subsubsection{7.2 DAG Affinity
Scheduling}\label{dag-affinity-scheduling}}

\begin{longtable}[]{@{}lll@{}}
\toprule\noalign{}
Strategy & Description & Best For \\
\midrule\noalign{}
\endhead
\bottomrule\noalign{}
\endlastfoot
\texttt{weight\_based} & Assigns to lowest-load worker & General use \\
\texttt{round\_robin} & Even distribution & Uniform workloads \\
\texttt{least\_loaded} & Real-time CPU metrics & Dynamic workloads \\
\texttt{pinned} & Fixed worker assignment & Isolation \\
\end{longtable}

\hypertarget{fault-tolerance}{%
\subsubsection{7.3 Fault Tolerance}\label{fault-tolerance}}

\begin{verbatim}
┌─────────────────────────────────────────────────────────────────────┐
│                     Fault Tolerance Mechanism                        │
├─────────────────────────────────────────────────────────────────────┤
│                                                                      │
│  Health Monitor Loop (every 2 seconds):                              │
│  ┌─────────────────────────────────────────────────────────────────┐│
│  │  for each worker in workers:                                    ││
│  │      if worker.missed_heartbeats > 3:                           ││
│  │          mark_unhealthy(worker)                                 ││
│  │          migrate_dags(worker → healthy_workers)                 ││
│  │          attempt_restart(worker, backoff=exponential)           ││
│  └─────────────────────────────────────────────────────────────────┘│
│                                                                      │
│  Recovery Timeline:                                                  │
│  t=0s    Worker crash detected                                       │
│  t=6s    Third missed heartbeat → UNHEALTHY                          │
│  t=6s    DAGs migrated to healthy workers                            │
│  t=6s    Restart attempt #1 (backoff: 5s)                            │
│  t=11s   Restart attempt #2 (backoff: 10s)                           │
│  ...     Max backoff: 60s                                            │
└─────────────────────────────────────────────────────────────────────┘
\end{verbatim}

\begin{center}\rule{0.5\linewidth}{0.5pt}\end{center}

\hypertarget{a.-enterprise-operational-model}{%
\subsection{7A. Enterprise Operational
Model}\label{a.-enterprise-operational-model}}

A dataflow engine is only useful in production if it can be deployed,
secured, made highly available, and operated. DishtaYantra layers four
operational subsystems over the compute core, each pluggable and
configuration-driven.

\hypertarget{a.1-pluggable-storage-abstraction}{%
\subsubsection{7A.1 Pluggable Storage
Abstraction}\label{a.1-pluggable-storage-abstraction}}

DAG definitions, holiday calendars, and other persisted artifacts are
read and written through a \texttt{StorageProvider} interface rather
than direct filesystem calls. Four providers ship: local filesystem
(default), Amazon S3, Azure Blob Storage, and Google Cloud Storage.
Switching providers is purely a configuration change; the compute core
is unaware of where its definitions live, which enables shared
definitions across a horizontally-scaled fleet.

\hypertarget{a.2-authentication-and-authorization}{%
\subsubsection{7A.2 Authentication and
Authorization}\label{a.2-authentication-and-authorization}}

Users, roles, and API keys are stored in a relational database (SQLite
by default, PostgreSQL optionally) via a DAO layer. Passwords are hashed
with PBKDF2-SHA256; clear text is never persisted. Interactive sessions
use signed cookies; programmatic clients authenticate with API keys.
Role checks (\texttt{admin} versus standard) gate administrative
operations. A legacy flat-file user store is migrated once into the
database and then retired.

\hypertarget{a.3-high-availability}{%
\subsubsection{7A.3 High Availability}\label{a.3-high-availability}}

A configurable HA manager elects a single PRIMARY instance across a
fleet, with automatic failover. Four election backends are provided ---
ZooKeeper (znode election), Redis (lease key with TTL), an S3-object
lease, and an exclusive-socket mechanism requiring no external
dependency. On demotion an instance suspends its DAGs; on election it
resumes them. The active role is surfaced in the operator console on
every page.

\begin{verbatim}
┌──────────────┐   election    ┌──────────────┐
│  Instance A  │◄────────────► │  Instance B  │
│  PRIMARY     │   (zk/redis/  │  SECONDARY   │
│  runs DAGs   │    s3/socket) │  DAGs paused │
└──────┬───────┘               └──────┬───────┘
       │  failure                     │ promote on
       └──────────────────────────────┘ PRIMARY loss
\end{verbatim}

\hypertarget{a.4-market-aware-scheduling}{%
\subsubsection{7A.4 Market-Aware
Scheduling}\label{a.4-market-aware-scheduling}}

Beyond simple time windows, each DAG may carry a schedule comprising a
start time plus a duration (e.g.~\texttt{09:30} + \texttt{6h30m}), a
day-of-week allow/deny list, and one or more holiday calendars
(e.g.~USA, Canada market calendars). A DAG is active only when the time
window matches, the weekday is permitted, and the date is not a holiday
in any followed calendar. Time windows that wrap past midnight are
handled correctly. Schedule edits --- to a DAG's config or to the
calendar files --- are honored by the running system within five
minutes, without a restart.

\hypertarget{a.5-format-agnostic-configuration}{%
\subsubsection{7A.5 Format-Agnostic
Configuration}\label{a.5-format-agnostic-configuration}}

System configuration is read from either YAML or Java-style properties;
the loader auto-detects the format and resolves
\texttt{\$\{VAR:default\}} placeholders against other keys, environment
variables, and command-line overrides, in that precedence order. Missing
required properties cause a fail-fast startup error that names the exact
missing key rather than silently defaulting.

\begin{center}\rule{0.5\linewidth}{0.5pt}\end{center}

\hypertarget{b.-recent-advances-columnar-transport-and-headless-orchestration}{%
\subsection{7B. Recent Advances: Columnar Transport and Headless
Orchestration}\label{b.-recent-advances-columnar-transport-and-headless-orchestration}}

The architecture described above has been extended with two capabilities
that preserve its defining invariant --- recomputation only when a
node's input \emph{actually changes} --- while broadening the
performance envelope and the operational model.

\hypertarget{b.1-arrow-native-columnar-transport}{%
\subsubsection{7B.1 Arrow-Native Columnar
Transport}\label{b.1-arrow-native-columnar-transport}}

The default edge transport carries one logical message as a Python
dictionary wrapped in a small envelope. For high-volume, uniformly-typed
streams this imposes a per-message overhead that dominates once the
calculator itself is cheap. We add an opt-in transport in which an edge
carries an immutable Apache Arrow \texttt{RecordBatch} and passes it
\emph{by reference} between nodes, eliminating per-row Python object
churn and copies.

Crucially, the equality gate is retained verbatim: because a
\texttt{RecordBatch} is immutable, the gate compares successive batches
with \texttt{RecordBatch.equals}, so a node still recomputes only when
its columnar input differs. Value handling is routed through a single
dispatch layer that treats the dictionary path byte-for-byte as before
and adds a batch path alongside it, so the change is additive and the
existing engine semantics are provably unchanged. A projection
transformer provides zero-copy column select/rename for the common
``narrow the wire'' case.

On a representative pipeline this columnar path is
\textbf{\textasciitilde2.29× faster end-to-end than the
dictionary-envelope path while producing byte-identical output} (zero
mismatches). The isolated kernel speedup is far larger
(\textasciitilde11.8×); the remaining gap is per-tick node-loop overhead
rather than transport, which usefully localizes where further work pays
off. Because the batches are standard Arrow, this transport also
establishes the substrate for zero-copy polyglot hand-off via the Arrow
C Data Interface to C++, Rust, and JVM calculators.

\hypertarget{b.2-headless-execution-and-orchestration}{%
\subsubsection{7B.2 Headless Execution and
Orchestration}\label{b.2-headless-execution-and-orchestration}}

The same DAG definitions that run behind the web service can also run
\emph{headless}: a runner starts the compute graph without the webapp,
optionally replays a bounded input feed, and treats completion as either
a target record count or quiescence (no dirty nodes and an empty input
queue across several polls). It drains in-flight work, writes a
machine-readable summary, and exits with a conventional status code ---
making a dataflow a first-class batch job.

Above the runner sits a lightweight control-plane/worker dispatch layer.
A dispatch calculator launches worker jobs that are \textbf{idempotent}
(duplicate keys are de-duplicated, so replays are safe),
\textbf{asynchronous} (workers run as child processes with a supervising
waiter), \textbf{bounded} (a configurable concurrency limit with a
pending pool prevents resource exhaustion), and \textbf{observable}
(each worker's summary is read back and can be republished downstream).
This gives backfills and scheduled batch runs the same dataflow
semantics --- including the change-only gate --- as the live service,
without bespoke orchestration code.

\begin{center}\rule{0.5\linewidth}{0.5pt}\end{center}

\hypertarget{related-work}{%
\subsection{8. Related Work}\label{related-work}}

\hypertarget{stream-processing-frameworks}{%
\subsubsection{8.1 Stream Processing
Frameworks}\label{stream-processing-frameworks}}

\textbf{Apache Flink} {[}2{]} provides exactly-once stream processing
with sophisticated windowing semantics. However, Flink's Python API
(PyFlink) introduces significant overhead due to cross-language
serialization.

\textbf{Apache Spark} {[}3{]} pioneered in-memory distributed computing
with the RDD abstraction. Spark Streaming processes data in
micro-batches, introducing minimum latencies of 100ms-1s.

\textbf{Naiad} {[}4{]} introduced timely dataflow for iterative
computations with low-latency incremental updates.

\hypertarget{zero-copy-technologies}{%
\subsubsection{8.2 Zero-Copy
Technologies}\label{zero-copy-technologies}}

\textbf{Apache Arrow} {[}12{]} provides a columnar memory format
enabling zero-copy data sharing. DishtaYantra extends this concept to
arbitrary data structures via LMDB's memory-mapped B-tree.

\textbf{LMDB} {[}8{]} is a high-performance key-value store used by
DishtaYantra for zero-copy data exchange.

\hypertarget{workflow-orchestration}{%
\subsubsection{8.3 Workflow
Orchestration}\label{workflow-orchestration}}

\textbf{Prefect} {[}13{]} and \textbf{Dagster} {[}14{]} are modern
Python-native workflow orchestration tools, but lack DishtaYantra's
focus on microsecond-latency processing.

\begin{center}\rule{0.5\linewidth}{0.5pt}\end{center}

\hypertarget{experimental-evaluation}{%
\subsection{9. Experimental Evaluation}\label{experimental-evaluation}}

\hypertarget{experimental-setup}{%
\subsubsection{9.1 Experimental Setup}\label{experimental-setup}}

\textbf{Hardware:} - CPU: AMD Ryzen 9 5900X (12-core, 24-thread) @
3.7GHz - RAM: 64GB DDR4-3200 - Storage: NVMe SSD (Samsung 980 PRO)

\textbf{Software:} - OS: Ubuntu 24.04 LTS - Python: 3.12.1 - Java:
OpenJDK 21 - C++ Compiler: GCC 13.2 - Rust: 1.75.0

\hypertarget{workload-descriptions}{%
\subsubsection{9.2 Workload Descriptions}\label{workload-descriptions}}

We evaluate DishtaYantra using two representative workloads:

\begin{longtable}[]{@{}
  >{\raggedright\arraybackslash}p{(\columnwidth - 4\tabcolsep) * \real{0.2500}}
  >{\raggedright\arraybackslash}p{(\columnwidth - 4\tabcolsep) * \real{0.3250}}
  >{\raggedright\arraybackslash}p{(\columnwidth - 4\tabcolsep) * \real{0.4250}}@{}}
\toprule\noalign{}
\begin{minipage}[b]{\linewidth}\raggedright
Workload
\end{minipage} & \begin{minipage}[b]{\linewidth}\raggedright
Description
\end{minipage} & \begin{minipage}[b]{\linewidth}\raggedright
Characteristics
\end{minipage} \\
\midrule\noalign{}
\endhead
\bottomrule\noalign{}
\endlastfoot
\textbf{W1: Passthrough} & Data passes through DAG without computation &
Measures framework overhead, I/O latency, serialization cost. Represents
best-case latency and throughput achievable by the system. \\
\textbf{W2: Mathematical Aggregation} & Statistical calculations (mean,
stddev, min, max, percentiles) on numerical arrays & CPU-bound workload
testing calculator efficiency. Represents typical compute-bound
analytics scenario. \\
\end{longtable}

\hypertarget{latency-results}{%
\subsubsection{9.3 Latency Results}\label{latency-results}}

\begin{longtable}[]{@{}lllll@{}}
\toprule\noalign{}
Payload Size & P50 & P95 & P99 & Max \\
\midrule\noalign{}
\endhead
\bottomrule\noalign{}
\endlastfoot
1 KB & 45 μs & 78 μs & 120 μs & 250 μs \\
10 KB & 65 μs & 110 μs & 180 μs & 350 μs \\
100 KB & 180 μs & 320 μs & 480 μs & 850 μs \\
1 MB & 850 μs & 1.2 ms & 1.8 ms & 3.2 ms \\
\end{longtable}

\emph{Table 1: End-to-end latency percentiles for W1 (Passthrough)
workload}

\hypertarget{throughput-results}{%
\subsubsection{9.4 Throughput Results}\label{throughput-results}}

\begin{longtable}[]{@{}lll@{}}
\toprule\noalign{}
Workers & W1 (Passthrough) & W2 (Math Aggregation) \\
\midrule\noalign{}
\endhead
\bottomrule\noalign{}
\endlastfoot
1 & 180,000 msg/s & 95,000 msg/s \\
2 & 350,000 msg/s & 185,000 msg/s \\
4 & 680,000 msg/s & 360,000 msg/s \\
6 & 1,000,000 msg/s & 530,000 msg/s \\
8 & 1,320,000 msg/s & 700,000 msg/s \\
10 & 1,620,000 msg/s & 860,000 msg/s \\
12 & 1,950,000 msg/s & 1,030,000 msg/s \\
\end{longtable}

\emph{Table 2: Maximum sustainable throughput (messages/second) by
worker count}

\hypertarget{scalability-analysis}{%
\subsubsection{9.5 Scalability Analysis}\label{scalability-analysis}}

\begin{verbatim}
              Throughput Scalability (W1 Workload)
              
  Throughput   │
  (msg/s)      │
               │
  2.0M ┤                                        ★ Actual
       │                                    ★
  1.5M ┤                                ★
       │                            ★   ---- Ideal Linear
  1.0M ┤                        ★---
       │                    ★---
  0.5M ┤                ★---
       │            ★---
       │        ★---
    0  ┼────★───┬───┬───┬───┬───┬───┬───┬───┬───┬───┬───┬
       0   1   2   3   4   5   6   7   8   9  10  11  12
                         Number of Workers
                         
  Parallel Efficiency at 12 workers: 90%
\end{verbatim}

\emph{Figure 5: Throughput scalability showing near-linear scaling up to
12 workers}

\hypertarget{lmdb-zero-copy-impact}{%
\subsubsection{9.6 LMDB Zero-Copy Impact}\label{lmdb-zero-copy-impact}}

\begin{longtable}[]{@{}llll@{}}
\toprule\noalign{}
Scenario & Without LMDB & With LMDB & Improvement \\
\midrule\noalign{}
\endhead
\bottomrule\noalign{}
\endlastfoot
100KB to C++ & 5.2 ms & 52 μs & 100× \\
1MB to Rust & 48 ms & 195 μs & 246× \\
10MB to Java & 520 ms & 2.1 ms & 248× \\
\end{longtable}

\emph{Table 3: Impact of LMDB zero-copy on cross-language communication}

\begin{center}\rule{0.5\linewidth}{0.5pt}\end{center}

\hypertarget{discussion}{%
\subsection{10. Discussion}\label{discussion}}

\hypertarget{design-trade-offs}{%
\subsubsection{10.1 Design Trade-offs}\label{design-trade-offs}}

\textbf{Memory vs.~Latency:} LMDB's memory-mapped design trades memory
usage for latency. The default 1GB map size accommodates most workloads.

\textbf{Process Isolation vs.~Communication Overhead:} The
multiprocessing architecture eliminates GIL contention but introduces
IPC overhead. LMDB pub/sub minimizes this cost.

\textbf{Flexibility vs.~Complexity:} Supporting five languages increases
implementation complexity but enables optimal algorithm selection.

\hypertarget{limitations}{%
\subsubsection{10.2 Limitations}\label{limitations}}

\begin{enumerate}
\def\labelenumi{\arabic{enumi}.}
\tightlist
\item
  \textbf{Single-Node Focus:} Current implementation targets single-node
  deployments
\item
  \textbf{JVM Startup Time:} Java calculators require JVM initialization
  (\textasciitilde2-5 seconds)
\item
  \textbf{Memory Footprint:} Each worker process maintains independent
  memory
\end{enumerate}

\hypertarget{future-work}{%
\subsubsection{10.3 Future Work}\label{future-work}}

\begin{enumerate}
\def\labelenumi{\arabic{enumi}.}
\tightlist
\item
  \textbf{Distributed Mode:} Multi-node deployment with automatic DAG
  partitioning
\item
  \textbf{GPU Calculators:} CUDA/ROCm integration for ML inference
\item
  \textbf{Checkpoint/Recovery:} Exactly-once semantics with distributed
  checkpointing
\item
  \textbf{WebAssembly Calculators:} Sandboxed calculator execution via
  WASM
\end{enumerate}

\begin{center}\rule{0.5\linewidth}{0.5pt}\end{center}

\hypertarget{conclusion}{%
\subsection{11. Conclusion}\label{conclusion}}

We presented \textbf{DishtaYantra}, a high-performance multi-language
dataflow framework that addresses critical limitations in existing
systems. Through LMDB-based zero-copy data exchange, unified
multi-language calculator support, and GIL-free worker parallelism,
DishtaYantra achieves:

\begin{itemize}
\tightlist
\item
  \textbf{100-250× latency improvements} over traditional JSON
  serialization for large payloads
\item
  \textbf{Sub-millisecond P99 latency} for payloads up to 100KB
\item
  \textbf{90\% parallel efficiency} at 12 workers on consumer hardware
\item
  \textbf{Unified support} for Python, Java, C++, Rust, and REST
  calculators
\item
  \textbf{Comprehensive pub/sub} integration with 10+ message brokers
\end{itemize}

DishtaYantra is open-source and available at
https://github.com/ajsinha/dishtayantra.

\begin{center}\rule{0.5\linewidth}{0.5pt}\end{center}

\hypertarget{references}{%
\subsection{References}\label{references}}

{[}1{]} T. Akidau et al., ``MillWheel: Fault-tolerant stream processing
at internet scale,'' VLDB Endowment, 2013.

{[}2{]} P. Carbone et al., ``Apache Flink: Stream and batch processing
in a single engine,'' IEEE Data Engineering Bulletin, 2015.

{[}3{]} M. Zaharia et al., ``Apache Spark: A unified engine for big data
processing,'' Communications of the ACM, 2016.

{[}4{]} D. G. Murray et al., ``Naiad: A timely dataflow system,'' SOSP,
2013.

{[}5{]} M. Kleppmann, ``Designing Data-Intensive Applications,''
O'Reilly Media, 2017.

{[}6{]} D. Beazley, ``Understanding the Python GIL,'' PyCon, 2010.

{[}7{]} W. M. Johnston et al., ``Advances in dataflow programming
languages,'' ACM Computing Surveys, 2004.

{[}8{]} H. Chu, ``LMDB: Lightning Memory-Mapped Database,''
https://www.symas.com/lmdb, 2011.

{[}9{]} W. Jakob et al., ``pybind11,''
https://github.com/pybind/pybind11, 2016.

{[}10{]} PyO3 Contributors, ``PyO3: Rust bindings for Python,''
https://pyo3.rs/, 2017.

{[}11{]} B. Dagenais, ``Py4J,'' https://www.py4j.org/, 2009.

{[}12{]} Apache Software Foundation, ``Apache Arrow,''
https://arrow.apache.org/, 2016.

{[}13{]} Prefect Technologies, ``Prefect,'' https://www.prefect.io/,
2018.

{[}14{]} Elementl, ``Dagster,'' https://dagster.io/, 2019.

\begin{center}\rule{0.5\linewidth}{0.5pt}\end{center}

\hypertarget{disclaimer-and-legal-notice}{%
\subsection{Disclaimer and Legal
Notice}\label{disclaimer-and-legal-notice}}

\textbf{DISCLAIMER OF WARRANTIES}

THE SOFTWARE IS PROVIDED ``AS IS'', WITHOUT WARRANTY OF ANY KIND,
EXPRESS OR IMPLIED, INCLUDING BUT NOT LIMITED TO THE WARRANTIES OF
MERCHANTABILITY, FITNESS FOR A PARTICULAR PURPOSE AND NONINFRINGEMENT.
IN NO EVENT SHALL THE AUTHORS OR COPYRIGHT HOLDERS BE LIABLE FOR ANY
CLAIM, DAMAGES OR OTHER LIABILITY, WHETHER IN AN ACTION OF CONTRACT,
TORT OR OTHERWISE, ARISING FROM, OUT OF OR IN CONNECTION WITH THE
SOFTWARE OR THE USE OR OTHER DEALINGS IN THE SOFTWARE.

\textbf{LIMITATION OF LIABILITY}

TO THE MAXIMUM EXTENT PERMITTED BY APPLICABLE LAW, IN NO EVENT SHALL THE
AUTHOR, CONTRIBUTORS, OR COPYRIGHT HOLDERS BE LIABLE FOR ANY DIRECT,
INDIRECT, INCIDENTAL, SPECIAL, EXEMPLARY, OR CONSEQUENTIAL DAMAGES
(INCLUDING, BUT NOT LIMITED TO, PROCUREMENT OF SUBSTITUTE GOODS OR
SERVICES; LOSS OF USE, DATA, OR PROFITS; OR BUSINESS INTERRUPTION)
HOWEVER CAUSED AND ON ANY THEORY OF LIABILITY, WHETHER IN CONTRACT,
STRICT LIABILITY, OR TORT (INCLUDING NEGLIGENCE OR OTHERWISE) ARISING IN
ANY WAY OUT OF THE USE OF THIS SOFTWARE, EVEN IF ADVISED OF THE
POSSIBILITY OF SUCH DAMAGE.

\textbf{USAGE RESPONSIBILITY}

Users of DishtaYantra are solely responsible for: - Ensuring compliance
with all applicable laws and regulations in their jurisdiction -
Validating the software's suitability for their specific use case -
Implementing appropriate security measures for production deployments -
Backing up data and maintaining disaster recovery procedures - Testing
thoroughly in non-production environments before deployment

\textbf{NO PROFESSIONAL ADVICE}

This software and documentation do not constitute professional advice of
any kind. Users should consult with qualified professionals for legal,
financial, security, or other specialized guidance.

\textbf{THIRD-PARTY COMPONENTS}

DishtaYantra incorporates third-party open-source components, each
subject to their respective licenses. Users are responsible for
compliance with all applicable third-party licenses.

\textbf{INDEMNIFICATION}

By using this software, you agree to indemnify and hold harmless the
author and contributors from any claims, damages, or expenses arising
from your use of the software.

\begin{center}\rule{0.5\linewidth}{0.5pt}\end{center}

\emph{© 2025 Ashutosh Sinha. All rights reserved.}

\emph{Contact: ajsinha@gmail.com}

\emph{Repository: https://github.com/ajsinha/dishtayantra}

\end{document}
